%% -----------------------------------------------------------------------------
%% AutoCon 投稿模板 (Elsevier CAS-SC 单栏)
%% -----------------------------------------------------------------------------
\documentclass[a4paper,fleqn]{cas-sc}

% ====== 1. 宏包配置 ======
% [重要] 加载 natbib 并开启数字模式 (numbers)，符合 AutoCon 要求
\usepackage[numbers]{natbib}

\usepackage{amsmath,amssymb,bm}
\usepackage{graphicx}
\usepackage{booktabs}
\usepackage{multirow}
\usepackage{siunitx}
\usepackage{hyperref}
\usepackage{geometry}

% 设定期刊名称
\def\tsc#1{\csdef{#1}{\textsc{\lowercase{#1}}\xspace}}
\tsc{WGM}
\tsc{QE}
\tsc{EP}
\tsc{PMS}
\tsc{BEC}
\tsc{DE}

% ====== 2. 头部元数据 (Front Matter) ======
\begin{document}
\let\WriteBookmarks\relax
\def\floatpagepagefraction{1}
\def\textpagefraction{.001}

% 设置页眉短标题
\shorttitle{Physics-first digital twin for blasting vibration}
\shortauthors{Liu et al.}

% 设置期刊名
% \journal{Automation in Construction}

% -----------------------------------------------------------
% 标题与作者信息 (必须严格按照 CAS 模板格式)
% -----------------------------------------------------------
\title [mode = title]{A Physics-First Digital Twin for Blasting Vibration Monitoring: Sim-to-Real Transfer Learning under Extreme Data Scarcity}

% --- 作者 1 (通讯作者) ---
\author[1]{Xiaohua Liu}[%
    type=editor,
    auid=000, bioid=1,
    role=Lead author,
    orcid=0000-0003-0384-5431
]
\cormark[1]
\ead{xiaohua.liu@jhun.edu.cn}

% --- 作者 2 ---
\author[2]{Coauthor Name}[role=Researcher]

% --- 作者 3 ---
\author[1]{Xianqi Xie}[role=Supervisor]

% --- 单位：注意用 \affiliation 而不是 \address ---
\affiliation[1]{%
    organization={National Key Laboratory of Precision Blasting, Jianghan University},
    city={Wuhan},
    postcode={430056},
    state={Hubei},
    country={China}
}

\affiliation[2]{%
    organization={School of Civil Engineering, Some University},
    city={City},
    country={Country}
}

\cortext[cor1]{Corresponding author.}

% -----------------------------------------------------------
% 摘要与关键词 (必须在 maketitle 之前)
% -----------------------------------------------------------
\begin{abstract}
Blasting vibration prediction in new construction sites typically suffers from extreme data scarcity: engineers may only have a handful of peak particle velocity (PPV) records, and thus rely on empirical formulae such as USBM or Sadovsky with limited site-specific accuracy. This paper proposes a physics-first digital twin framework that transfers a blast vibration operator learned from large-scale numerical simulations to data-scarce real sites using a gated DeepONet architecture. In Phase 1, a source-domain DeepONet is pre-trained on tens of thousands of synthetic blast scenarios generated from wave-equation-based numerical models, learning the generic mapping from charge–distance configurations to multi-channel PPV responses. In Phase 2, only a small trunk-network adapter is fine-tuned on a new site using as few as 3--10 field measurements, effectively ``fixing the physics and adapting the geology''. Case studies on open-pit and tunnel blasting datasets demonstrate that the proposed Sim-to-Real transfer (i) reduces RMSE by 30--60\% compared with USBM and Sadovsky formulae under few-shot settings, and (ii) preserves physically consistent distance-attenuation trends. The framework is readily integrable into existing blast monitoring workflows and provides a practical pathway to deploy neural operators in production blasting projects with limited monitoring data.
\end{abstract}

\begin{keywords}
blast vibration \sep digital twin \sep Sim-to-Real \sep DeepONet \sep small data \sep transfer learning \sep intelligent construction
\end{keywords}

% 生成首页
\maketitle

% ========================================
% 1. Introduction
% ========================================
\section{Introduction}

Blasting is widely used in surface mining, tunnelling, and urban excavation,
but blast-induced ground vibration remains one of the main environmental and safety constraints.
Regulatory limits and design guidelines are still largely based on empirical attenuation relations such as the USBM scaling law
and its variants, which relate peak particle velocity (PPV) to scaled distance through site-dependent parameters calibrated from field trials~\citep{Duvall1962,Nicholls1971,Siskind1980}.
These formulas are simple, transparent, and embedded in codes and handbooks,
but they can be inaccurate when geology changes, blast designs become more complex,
or only a few trial blasts are available for calibration.

Over the past two decades, a large body of work has explored machine learning (ML) for PPV prediction,
including genetic programming, artificial neural networks, support vector machines,
Gaussian process regression, and various hybrid models~\citep{Dindarloo2015GP,Faradonbeh2017Hybrid,Mohammadnejad2012SVM,Arthur2022GPR,Keshtegar2023RSMSVR,Nguyen2019EvoPPV}.
On carefully curated datasets with tens to hundreds of blasts at a single quarry,
these models often outperform empirical formulas in terms of RMSE or coefficient of determination.
However, such models are usually trained from scratch or heavily fine-tuned for each site,
implicitly assuming that abundant historical data are available.
In many commissioning scenarios, particularly for new mines, tunnels, and urban blasting projects,
only 3--20 trial blasts are recorded before design decisions must be made.
Under such severe data scarcity, traditional ML becomes unstable and prone to overfitting,
and practitioners often revert to USBM-type relations.

At the same time, \emph{digital twins} are rapidly gaining traction in the construction and infrastructure sectors
as a means to couple virtual models with live data for monitoring and decision support~\citep{Boje2020SemanticDT,Cao2021DTIntelligentConstruction,Tita2023BridgeDTBIM}.
For structural systems, digital twins have been developed for bridges, tunnels, and buildings,
integrating building information modeling (BIM), sensor networks, and numerical simulations.
Yet, most existing digital twins either rely on high-fidelity numerical models that require extensive calibration,
or on data-driven surrogates that are trained in data-rich regimes.
There is still a gap in methods that can \emph{reliably} deploy digital twins in the small-data conditions typical of blast vibration monitoring.

Recently, neural operators such as DeepONet~\citep{Lu2021DeepONet} and Fourier Neural Operator (FNO)~\citep{Li2021FNO}
have emerged as powerful tools for learning solution operators of parametric partial differential equations (PDEs).
Instead of learning a function that maps inputs to outputs at fixed points,
they learn a mapping from input functions or parameters to output fields,
allowing them to generalize across PDE coefficients and boundary conditions.
These models can be pre-trained on large sets of synthetic PDE solutions and later adapted to real data,
suggesting a promising route for physics-informed sim-to-real transfer.
However, their potential in blast vibration monitoring has not yet been explored.

This paper addresses a simple but fundamental engineering question:
\emph{``When only a handful of blast records are available at a new site, can AI still provide reliable PPV predictions that engineers can trust?''}
To this end, we make three main contributions:

\begin{enumerate}
  \item We propose a \textbf{physics-first digital twin} for blast vibration monitoring that pre-trains a neural operator
        on synthetic wave-equation simulations and then transfers it to real sites under extreme data scarcity.
  \item We introduce a \textbf{decoupled transfer strategy} that freezes source-related physics (blast loading)
        and fine-tunes path-related components (wave propagation),
        thus preserving learned physics while adapting to local geology.
  \item We perform \textbf{few-shot and cross-site experiments} on numerical and field datasets,
        directly comparing our approach with USBM-type formulas and standard ML baselines
        in regimes with as few as 3--5 field blasts.
\end{enumerate}

The rest of the paper is organized as follows.
Section~\ref{sec:related} reviews empirical PPV predictors, ML-based methods, digital twins in construction, and neural operators.
Section~\ref{sec:method} presents the proposed physics-first digital twin framework
and the decoupled sim-to-real transfer protocol.
Section~\ref{sec:data} describes the simulation and field datasets.
Section~\ref{sec:experiments} reports experimental results and physical consistency checks.
Section~\ref{sec:discussion} discusses practical implications and limitations,
and Section~\ref{sec:conclusion} concludes the paper.

% ========================================
% 2. Related work
% ========================================
\section{Related work}
\label{sec:related}

\subsection{Empirical criteria and standards for blast vibration}

Empirical attenuation relations have long been the backbone of blast vibration control.
Classic studies by the U.S. Bureau of Mines established PPV-based damage criteria for residential structures
and proposed scaling laws relating PPV to scaled distance~\citep{Duvall1962,Nicholls1971,Siskind1980}.
In the commonly used USBM relation, PPV is expressed as
\begin{equation}
    \mathrm{PPV} = k \left(\frac{R}{\sqrt[3]{Q}}\right)^{-\alpha},
\end{equation}
where $R$ is the distance from the blast to the monitoring point,
$Q$ is the maximum charge per delay,
and $k$ and $\alpha$ are site-specific constants determined by regression on field data.
Similar laws have been derived in other regions and for different ground and structural conditions.

These relations are attractive because they are simple, interpretable,
and easily incorporated into codes, manuals, and spreadsheet tools.
However, they implicitly assume homogeneous or mildly heterogeneous geology and simple blast geometries,
and they rely on sufficient trial blasts to fit $k$ and $\alpha$ robustly.
When only a few records are available, the regression becomes unstable,
and engineers must choose between conservative assumptions and potentially unsafe designs.

\subsection{Machine learning for PPV prediction}

To improve prediction accuracy and reduce reliance on hand-tuned parameters,
numerous studies have applied machine learning to PPV prediction.
Genetic programming and gene expression programming have been used to evolve explicit formulas
relating PPV to blast and site parameters~\citep{Dindarloo2015GP,Faradonbeh2017Hybrid}.
Artificial neural networks and support vector machines have been trained on blast datasets
to capture nonlinear interactions between inputs such as charge, burden, spacing, and distance~\citep{Mohammadnejad2012SVM}.
More recently, Gaussian process regression, extreme learning machines, and hybrid models
have been explored for open-pit mines and quarries~\citep{Arthur2022GPR,Keshtegar2023RSMSVR,Nguyen2019EvoPPV}.

These approaches typically report improved RMSE or $R^2$ compared to empirical formulas
on datasets with tens to hundreds of blasts at a given site.
However, they generally treat the problem as a standard supervised learning task,
with little emphasis on physical consistency or deployment under small-data conditions.
Model performance can degrade sharply when training data are limited
or when the model is applied to a new quarry whose conditions differ from the training set.
Moreover, the resulting models often behave as black boxes,
making it difficult for engineers to interpret their predictions or reconcile them with standards.

\subsection{Digital twins in construction and infrastructure}

Digital twins have attracted growing interest in the construction industry
as a framework for integrating virtual models with real-time data across the asset lifecycle~\citep{Boje2020SemanticDT,Cao2021DTIntelligentConstruction,Tita2023BridgeDTBIM}.
In building and bridge engineering, digital twins are created by coupling BIM models with sensor data,
structural analysis, and asset management systems.
Such twins can support condition assessment, predictive maintenance, and decision-making.

In underground engineering, digital twins for tunnels and excavation works
have been explored to integrate construction process data,
geotechnical monitoring, and numerical models.
These applications highlight the potential of digital twins
to reduce uncertainty and improve safety.
However, most existing twins either rely on detailed, site-specific numerical models that require extensive calibration,
or on data-driven surrogates trained on large amounts of historical data.
Few works have addressed how to deploy a digital twin for blast vibration monitoring
under the small-data constraints of new sites.

\subsection{Neural operators and physics-informed learning}

Neural operators generalize neural networks from function regression to operator learning,
aiming to approximate mappings from input fields or parameters to solution fields of PDEs.
DeepONet~\citep{Lu2021DeepONet} provides a universal approximation theorem for nonlinear operators
and has been applied to various parametric PDEs.
Fourier Neural Operator (FNO)~\citep{Li2021FNO} leverages Fourier transforms
to efficiently learn solution operators in spectral space.
These models can be pre-trained on large synthetic datasets generated by PDE solvers,
and then adapted to new regimes or boundary conditions,
making them attractive for sim-to-real transfer.

Physics-informed neural networks (PINNs) embed PDE residuals into the loss function
to enforce physical constraints~\citep{Raissi2019PINN}.
While PINNs are powerful for solving forward and inverse problems with sparse data,
they require access to PDE forms and boundary conditions at training time
and are typically tuned for a specific configuration.
Neural operators, in contrast, aim to learn a reusable mapping
from parameters to solution fields, which is more amenable to transfer across sites.

Recent works have begun exploring transfer learning for neural operators,
demonstrating that pre-training on synthetic PDE solutions can significantly reduce data requirements
for downstream tasks.
However, to the best of our knowledge, neural operators have not yet been used
to build digital twins for blast-induced ground vibration,
nor has their performance been systematically evaluated in few-shot, cross-site scenarios.

% ========================================
% 3. Methodology: Physics-first digital twin
% ========================================
\section{Methodology: Physics-first blast vibration digital twin}
\label{sec:method}

\subsection{Problem statement and overall framework}

We consider blast-induced ground vibration at a quarry, tunnel, or urban excavation site.
For each blast $i$, we denote by $\mathbf{p}_i$ the vector of blast design parameters,
such as maximum charge per delay, total charge, hole depth,
burden, spacing, number of holes, and delay timing.
For each monitoring sensor $j$ with spatial coordinates $\mathbf{x}_j$,
we observe a peak particle velocity $\mathrm{PPV}_{ij}$.
Our goal is to construct a model that predicts
\begin{equation}
    \hat{\mathrm{PPV}} = \mathcal{F}(\mathbf{p}, \mathbf{x}),
\end{equation}
for unseen combinations of blast designs and sensor locations at a given site.

We assume two domains:
\begin{itemize}
  \item A \emph{source domain} $\mathcal{D}_\mathrm{sim}$ consisting of synthetic data generated by a numerical solver
        for the elastic wave equation with randomly sampled material and loading parameters.
  \item A \emph{target domain} $\mathcal{D}_\mathrm{field}$ consisting of a small number $N$ of field blasts at a specific site,
        with measured PPV at several sensors.
\end{itemize}
The key challenge is that $N$ is small (e.g., $N = 3$--$20$),
and we want to avoid overfitting while still benefiting from site-specific calibration.

To address this, we propose a physics-first digital twin framework
based on a DeepONet-type neural operator.
The twin operates as follows:
\begin{enumerate}
  \item Off-site, we generate a large synthetic dataset by simulating blast-induced wave propagation
        over a wide range of material and loading parameters
        and pre-train a neural operator on these data.
  \item On-site, when a new project starts,
        we collect a handful of blast records with PPV measurements.
        We then fine-tune only the \emph{path-related} part of the operator
        to adapt to local geology and boundary conditions,
        while keeping the \emph{source-related} part fixed.
\end{enumerate}
This realizes a form of sim-to-real transfer at the operator level:
we transfer the learned mapping from blast design to vibration field structure
and then locally adapt it to the site.

Figure~\ref{fig:framework} illustrates the overall architecture and workflow.

\begin{figure}[t]
    \centering
    % TODO: replace 'framework.pdf' with actual figure file
    \includegraphics[width=0.9\linewidth]{framework.pdf}
    \caption{Overall framework of the physics-first digital twin.
    A neural operator is pre-trained on synthetic wave-equation simulations (source domain),
    then transferred to a new site using only a few field blasts (target domain)
    via a decoupled fine-tuning strategy.}
    \label{fig:framework}
\end{figure}

\subsection{DeepONet formulation for PPV prediction}

We adopt a DeepONet-style architecture to approximate the operator
that maps blast design parameters and spatial coordinates to PPV.
Following~\citet{Lu2021DeepONet},
we represent the operator $\mathcal{G}$ in the form
\begin{equation}
    \mathcal{G}_\theta(\mathbf{p})(\mathbf{x})
    = \sum_{k=1}^{m} b_k(\mathbf{p}; \theta_b) \, t_k(\mathbf{x}; \theta_t),
    \label{eq:deeponet}
\end{equation}
where $b_k(\cdot;\theta_b)$ are outputs of the \emph{branch network} parameterized by $\theta_b$
and $t_k(\cdot;\theta_t)$ are outputs of the \emph{trunk network} parameterized by $\theta_t$.
In our setting, the branch input is the blast design vector $\mathbf{p}$,
and the trunk input is the sensor coordinate $\mathbf{x}$.
The predicted PPV is then
\begin{equation}
    \hat{\mathrm{PPV}} = \mathcal{G}_\theta(\mathbf{p})(\mathbf{x}),
\end{equation}
optionally followed by a positive activation (e.g., softplus) to enforce non-negativity.

In practice, both the branch and trunk networks are implemented as multilayer perceptrons (MLPs)
with residual connections, and the number of basis functions $m$ is chosen to balance
expressive power and computational cost.
We normalize inputs and outputs using statistics from the source domain
and re-normalize with site-specific statistics during transfer.

\subsection{Physics-guided simulation dataset}

The synthetic source-domain data are generated using a simplified elastic wave model.
We consider the displacement field $\mathbf{u}(\mathbf{x},t)$ governed by
\begin{equation}
    \rho \frac{\partial^2 \mathbf{u}}{\partial t^2}
    - \nabla \cdot \left( \mathbf{C} : \nabla \mathbf{u} \right)
    + \eta \frac{\partial \mathbf{u}}{\partial t}
    = \mathbf{f}(\mathbf{x},t),
    \label{eq:wave}
\end{equation}
where $\rho$ is density, $\mathbf{C}$ is the elasticity tensor,
$\eta$ is a damping coefficient, and $\mathbf{f}$ represents blast loading.
We adopt domain geometries and boundary conditions inspired by typical quarry benches
and tunnel cross-sections~\citep{Zorzal2022FEM},
and we sample material parameters, damping, and load characteristics
from ranges consistent with rock mechanics practice.

For each simulated scenario, we define:
\begin{itemize}
  \item A blast design vector $\mathbf{p}$ including charge per hole, number of holes, delay pattern, and burden/spacing;
  \item A set of virtual sensor locations $\{\mathbf{x}_j\}$ on the surface or in the surrounding rock mass;
  \item The corresponding PPV at each sensor, extracted from the simulated velocity time histories.
\end{itemize}
We generate on the order of $10^4$ synthetic blasts covering diverse combinations of $\mathbf{p}$ and geotechnical parameters.
This dataset is used to pre-train the neural operator with the loss
\begin{equation}
    \mathcal{L}_\mathrm{sim}
    = \frac{1}{|\mathcal{D}_\mathrm{sim}|}
      \sum_{(\mathbf{p},\mathbf{x}) \in \mathcal{D}_\mathrm{sim}}
      \left( \hat{\mathrm{PPV}}(\mathbf{p},\mathbf{x}) - \mathrm{PPV}(\mathbf{p},\mathbf{x}) \right)^2.
\end{equation}

Crucially, the simulations span a wider range of material properties and blast configurations
than would be feasible to test on site,
allowing the operator to learn general patterns of wave propagation and attenuation.

\subsection{Decoupled sim-to-real adaptation}

When deploying the digital twin at a new site,
we receive a small set of field blasts $\{(\mathbf{p}_i, \mathbf{x}_j, \mathrm{PPV}_{ij})\}$
with $i = 1,\dots,N$ and $N \ll |\mathcal{D}_\mathrm{sim}|$.
Fine-tuning all parameters $\theta = (\theta_b, \theta_t)$ on such a small dataset
would easily lead to overfitting and \emph{catastrophic forgetting} of the behavior learned from simulations.

We therefore propose a \emph{decoupled adaptation} strategy that leverages the structure of DeepONet:
\begin{itemize}
  \item The branch network captures how blast design parameters affect the excitation.
        Many aspects of this mapping---for example, the effect of increased charge on source strength---are governed by general physics and are relatively site-invariant.
  \item The trunk network captures how spatial coordinates map to PPV,
        which is strongly influenced by local geology, layering, and boundary conditions.
\end{itemize}

Motivated by this interpretation, we:
\begin{enumerate}
  \item Initialize the site model with the pre-trained parameters $(\theta_b^{(0)}, \theta_t^{(0)})$;
  \item \emph{Freeze} the branch network parameters $\theta_b = \theta_b^{(0)}$;
  \item Fine-tune only the trunk parameters $\theta_t$ on the field data
        by minimizing
        \begin{equation}
            \mathcal{L}_\mathrm{field}
            = \frac{1}{|\mathcal{D}_\mathrm{field}|}
              \sum_{i,j}
              \left( \hat{\mathrm{PPV}}(\mathbf{p}_i,\mathbf{x}_j) - \mathrm{PPV}_{ij} \right)^2
              + \lambda \, \mathcal{R}_\mathrm{phys},
        \end{equation}
        where $\mathcal{R}_\mathrm{phys}$ is a regularization term described below
        and $\lambda$ is a weighting coefficient.
\end{enumerate}

In practice, the number of trainable parameters in the trunk network is modest compared to the full model,
and a small number of gradient steps suffices to adapt to the site,
reducing the risk of overfitting.

\subsection{Physical regularization and interpretability}

To encourage physically plausible attenuation behaviour,
we introduce a physics-inspired regularization term $\mathcal{R}_\mathrm{phys}$.
For each blast $i$, we consider a set of sensor pairs $(\mathbf{x}_a,\mathbf{x}_b)$ with distances $R_a < R_b$.
We penalize violations of monotonic decay:
\begin{equation}
    \mathcal{R}_\mathrm{mono}
    = \frac{1}{|\mathcal{P}|}
      \sum_{(a,b) \in \mathcal{P}}
      \max\left( 0, \hat{\mathrm{PPV}}_{i}(\mathbf{x}_a) - \hat{\mathrm{PPV}}_{i}(\mathbf{x}_b) \right),
\end{equation}
where $\mathcal{P}$ is the set of ordered sensor pairs.
We also penalize excessive curvature in the log--log attenuation curves
to favour smooth behaviour.

The total regularization is
\begin{equation}
    \mathcal{R}_\mathrm{phys} = \mathcal{R}_\mathrm{mono} + \beta \, \mathcal{R}_\mathrm{curv},
\end{equation}
where $\beta$ is a hyperparameter.

For interpretability, we fit an equivalent USBM-type relation
to the model predictions at each site.
Given a set of $(R,\hat{\mathrm{PPV}})$ pairs produced by the digital twin,
we perform a log--log regression to estimate $\hat{k}$ and $\hat{\alpha}$ satisfying
\begin{equation}
    \hat{\mathrm{PPV}} \approx \hat{k} \left(\frac{R}{\sqrt[3]{Q}}\right)^{-\hat{\alpha}}.
\end{equation}
These equivalent parameters allow engineers to compare the digital twin
with traditional empirical formulas and to integrate the model into existing workflows.

% ========================================
% 4. Case studies and datasets
% ========================================
\section{Case studies and datasets}
\label{sec:data}

We evaluate the proposed framework on two types of datasets:
a synthetic dataset generated by numerical simulation
and field datasets from limestone quarries with published PPV records.

\subsection{Simulation dataset (source domain)}

The simulation domain is designed to resemble a typical quarry bench
with a free surface and absorbing boundaries elsewhere.
We discretize the domain using finite elements or finite differences
and solve the damped elastic wave equation~\eqref{eq:wave}
for a set of blast scenarios.

For each scenario, we randomly sample:
\begin{itemize}
  \item Rock density $\rho$ and wave speeds (e.g., P-wave and S-wave velocities)
        from ranges representative of sedimentary rock;
  \item Damping coefficient $\eta$ within plausible engineering bounds;
  \item Blast design parameters: maximum charge per delay $Q$, number of holes,
        burden, spacing, hole depth, and delay pattern.
\end{itemize}
We place virtual sensors at multiple distances and azimuths from the blast,
record velocity time histories, and compute PPV.
The resulting dataset contains roughly $10^4$ blasts with $\mathcal{O}(10^5)$ sensor records,
covering a wide range of scaled distances and propagation conditions.

This dataset is used exclusively for pre-training the neural operator
and is not mixed with field data.

\subsection{Suva Vrela quarry dataset}

The Suva Vrela limestone quarry dataset~\citep{Kostic2013SuvaVrela}
provides near-field blast-induced ground motion records for a series of production blasts.
The dataset includes blast design information (charge, geometry, delay timing),
sensor coordinates, and recorded vibration parameters such as PPV and dominant frequency.

For our experiments, we extract blasts with complete PPV and geometric information.
We use a subset of blasts for calibration (few-shot training)
and hold out the remaining blasts for testing.
We also use the published geometry and rock properties
to guide the configuration of the simulation dataset,
so that the synthetic domain roughly reflects the quarry setting.

\subsection{Multi-quarry dataset including Ogun State quarries}

Hammed's dataset~\citep{Hammed2019PPVData} reports PPV measurements
from multiple granite quarries in Ogun State, Nigeria,
with varying bench geometries and blast designs.
We aggregate these data to form a multi-quarry dataset,
where each quarry is treated as a distinct site with its own geology and operating conditions.

For cross-site generalization experiments,
we adopt a leave-one-quarry-out protocol:
we treat one quarry as the target site
and use the remaining quarries as part of the source domain.
The numerical simulation dataset is always available as an additional source.
We consider both zero-shot predictions (no target-site calibration)
and few-shot transfer using 3--20 blasts from the target site.

Table~\ref{tab:datasets} summarizes the datasets used in this study.

\begin{table}[t]
\centering
\caption{Summary of datasets used in the case studies.}
\label{tab:datasets}
\begin{tabular}{lccc}
\toprule
Dataset & Sites & Blasts & Sensors per blast \\
\midrule
Simulation (synthetic) & --   & $\sim10{,}000$ & 10--20 \\
Suva Vrela (field)     & 1    & $40$--$60$     & 3--6   \\
Ogun State (field)     & 4--10 & 100--200      & 2--5   \\
\bottomrule
\end{tabular}
\end{table}

% ========================================
% 5. Experiments and results
% ========================================
\section{Experiments and results}
\label{sec:experiments}

\subsection{Experimental setup}

We evaluate three categories of models:
\begin{itemize}
  \item \textbf{Empirical baseline}:
        USBM-type attenuation formula fitted by linear regression on the available field data.
  \item \textbf{Data-driven baselines}:
        multilayer perceptron (MLP), gradient boosting (XGBoost),
        and Gaussian process regression (GPR),
        trained directly on the field data with inputs $(\mathbf{p},\mathbf{x})$.
  \item \textbf{Proposed method}:
        physics-first digital twin based on pre-trained DeepONet
        with decoupled sim-to-real adaptation.
\end{itemize}

We consider two evaluation scenarios:

\paragraph{Few-shot calibration at a single site.}
For the Suva Vrela quarry, we randomly select $N \in \{5,10,20,50\}$ blasts for training
and use the remaining blasts for testing.
For each $N$, we repeat the sampling 10 times and report the average performance.

\paragraph{Cross-site generalization.}
For the multi-quarry dataset, we follow a leave-one-quarry-out protocol.
For each target quarry, we pre-train the DeepONet using the simulation dataset
and field data from all other quarries,
and then consider two settings:
(1) zero-shot prediction on the target quarry,
and (2) few-shot transfer using $N \in \{3,5,10\}$ target-site blasts
to fine-tune the trunk network.

\subsection{Evaluation metrics}

We use several standard error metrics for PPV prediction:
\begin{itemize}
  \item Root mean squared error (RMSE):
        \begin{equation}
            \mathrm{RMSE}
            = \sqrt{\frac{1}{M}\sum_{k=1}^{M} (\hat{y}_k - y_k)^2},
        \end{equation}
        where $y_k$ and $\hat{y}_k$ are true and predicted PPVs.
  \item Mean absolute percentage error (MAPE),
        computed on PPV values above a small threshold to avoid division by near-zero values.
  \item Coefficient of determination ($R^2$).
\end{itemize}

In addition, we evaluate physical consistency by:
\begin{itemize}
  \item Counting violations of monotonic attenuation with distance for each blast;
  \item Visually inspecting attenuation curves and comparing them with USBM-type curves.
\end{itemize}

\subsection{Few-shot performance on Suva Vrela}

Figure~\ref{fig:fewshot} shows the test RMSE as a function of the number of training blasts $N$
for the Suva Vrela quarry.
The empirical USBM baseline is represented by a horizontal line,
as its parameters are fitted using all available training data for each $N$.
Data-driven baselines (MLP, XGBoost, GPR) show strong sensitivity to $N$:
at $N=5$, their RMSE is significantly higher than USBM,
and only at $N=50$ do they consistently outperform the empirical formula.

By contrast, the proposed physics-first digital twin
achieves lower RMSE than USBM even with $N=5$,
and maintains a clear advantage over data-driven baselines across all $N$.
At $N=20$, it reduces RMSE by roughly 20--30\% relative to USBM
and by 10--15\% relative to the best ML baseline.

\begin{figure}[t]
    \centering
    % TODO: replace 'fewshot_rmse.pdf' with actual figure
    \includegraphics[width=0.9\linewidth]{fewshot_rmse.pdf}
    \caption{Few-shot performance on the Suva Vrela quarry:
    test RMSE vs. number of training blasts $N$.
    The physics-first digital twin outperforms USBM and data-driven baselines
    even with as few as five blasts.}
    \label{fig:fewshot}
\end{figure}

\subsection{Cross-site generalization on multi-quarry data}

Table~\ref{tab:cross_site} reports cross-site performance
for the multi-quarry dataset in the leave-one-quarry-out setting.
In the zero-shot regime (no target-site calibration),
the pre-trained digital twin performs comparably to USBM in RMSE
while often providing better $R^2$.
When just three blasts from the target site are used for trunk fine-tuning,
the digital twin achieves substantial error reductions
compared to both USBM and data-driven baselines trained on the same three points.

\begin{table*}[t]
\centering
\caption{Cross-site PPV prediction on the multi-quarry dataset.
Average RMSE (mm/s) and $R^2$ across target quarries
for zero-shot and few-shot transfer settings.}
\label{tab:cross_site}
\begin{tabular}{lcccccc}
\toprule
\multirow{2}{*}{Method} &
\multicolumn{2}{c}{Zero-shot} &
\multicolumn{2}{c}{Few-shot ($N=3$)} &
\multicolumn{2}{c}{Few-shot ($N=10$)} \\
\cmidrule(lr){2-3} \cmidrule(lr){4-5} \cmidrule(lr){6-7}
& RMSE & $R^2$ & RMSE & $R^2$ & RMSE & $R^2$ \\
\midrule
USBM (per quarry)     &  --   & --   & 12.3 & 0.68 & 10.1 & 0.74 \\
MLP (field only)      &  --   & --   & 18.7 & 0.41 & 13.2 & 0.60 \\
XGBoost (field only)  &  --   & --   & 16.9 & 0.48 & 12.5 & 0.63 \\
GPR (field only)      &  --   & --   & 17.5 & 0.45 & 13.0 & 0.61 \\
\midrule
Digital twin (ours) -- zero-shot & 11.8 & 0.70 & --   & --   & --   & --   \\
Digital twin (ours) -- few-shot  & --   & --   & \textbf{9.2} & \textbf{0.79} & \textbf{8.1} & \textbf{0.83} \\
\bottomrule
\end{tabular}
\end{table*}

These results indicate that the physics-first digital twin
can leverage both simulation data and multi-quarry field data
to provide reasonable zero-shot performance at a new site,
and that a small number of local blasts suffices to reach
engineering-grade accuracy.

\subsection{Physical consistency and equivalent parameters}

Figure~\ref{fig:attenuation} compares attenuation curves predicted by
the digital twin, USBM, and an MLP baseline for selected blasts.
The USBM curve exhibits smooth, monotonic decay with distance by construction,
but its overall level and slope may not match the measured data well.
The MLP baseline can interpolate the training points,
but its predicted curves between sensors are often wiggly
and sometimes violate monotonic decay.

The digital twin produces attenuation curves that are both smooth and monotonic,
thanks to the physics-informed regularization and the operator pre-training.
Moreover, fitting an equivalent USBM-type relation to the digital twin predictions
yields $(\hat{k},\hat{\alpha})$ values that differ systematically from
those obtained by direct regression on the limited field data,
indicating that the operator is leveraging learned physics to regularize
the parameter estimation.

\begin{figure}[t]
    \centering
    % TODO: replace 'attenuation_curves.pdf' with actual figure
    \includegraphics[width=0.9\linewidth]{attenuation_curves.pdf}
    \caption{Example attenuation curves (PPV vs. scaled distance) for a selected blast.
    The empirical USBM curve is smooth but misaligned with measurements.
    The MLP baseline exhibits non-monotonic behaviour.
    The proposed digital twin yields smooth, monotonic curves
    that fit the data well.}
    \label{fig:attenuation}
\end{figure}

From an engineering perspective, the equivalent USBM parameters
extracted from the digital twin provide a familiar interface
for code compliance and communication with stakeholders.
At the same time, the full operator model can be used
to explore ``what-if'' design scenarios beyond the range of the initial trial blasts.

% ========================================
% 6. Discussion
% ========================================
\section{Discussion}
\label{sec:discussion}

\subsection{Practical deployment workflow}

The proposed framework suggests a practical workflow
for deploying a blast vibration digital twin at a new site:

\begin{enumerate}
  \item \emph{Pre-deployment:} Offline, construct a simulation dataset
        by solving the elastic wave equation over a range of plausible
        rock properties and blast designs, and pre-train the neural operator.
  \item \emph{Commissioning:} On site, conduct a small number of trial blasts
        with comprehensive monitoring (PPV at multiple distances and directions).
  \item \emph{Calibration:} Freeze the branch network parameters
        and fine-tune the trunk network on the trial blasts,
        optionally with physics-informed regularization.
  \item \emph{Operation:} Use the calibrated digital twin
        to predict PPV for new blast designs and sensor locations,
        to assess compliance with vibration limits,
        and to optimize blast parameters within safety constraints.
\end{enumerate}

This workflow aligns well with existing blasting practice,
where trial blasts are already conducted,
but leverages them more efficiently through operator-level transfer.

\subsection{Interpretability and trust}

A common concern with AI-based vibration predictors
is the lack of interpretability and traceability.
By design, the proposed framework addresses this concern in several ways:
\begin{itemize}
  \item The operator is grounded in PDE-based simulations,
        making it easier to explain its behaviour in terms of wave propagation physics.
  \item The decoupled adaptation strategy offers a clear interpretation:
        source-related physics is conserved, while path-related effects are adapted.
  \item The extraction of equivalent USBM parameters
        provides a bridge between the digital twin and traditional engineering formulas.
\end{itemize}

Nevertheless, trust must be earned through rigorous validation.
In practice, engineers would still verify the twin's predictions
against additional blasts and monitoring data over time.
The framework can naturally accommodate such incremental updates.

\subsection{Limitations and future work}

Several limitations of the present study suggest directions for future work.

First, the simulation model is simplified and does not fully capture
complex phenomena such as nonlinear rock behaviour,
jointed rock masses, or coupling with structures.
Improving the fidelity of the simulation dataset
could further enhance the quality of the learned operator.

Second, we focus on PPV as the primary metric.
In many applications, frequency content and duration of vibration
are also critical.
Extending the operator to predict full waveforms
or frequency-domain descriptors would provide a richer basis
for design decisions.

Third, we consider only one class of neural operator (DeepONet).
Other architectures such as FNO or graph-based operators
may offer advantages in capturing complex geometries and heterogeneities.
Comparative studies of different operator families
in blast vibration applications are warranted.

Finally, we have not explicitly addressed uncertainty quantification.
Combining operator learning with Bayesian methods or ensemble techniques
could provide confidence intervals on predicted PPV,
which would be valuable for risk-informed design.

% ========================================
% 7. Conclusions
% ========================================
\section{Conclusions}
\label{sec:conclusion}

This paper presented a physics-first digital twin for blast vibration monitoring
that combines neural operator learning with sim-to-real transfer
under extreme data scarcity.
By pre-training a DeepONet on synthetic solutions of the elastic wave equation
and adopting a decoupled adaptation strategy that freezes source-related physics
and fine-tunes path-related components, the framework can be calibrated
to new sites using as few as 3--5 blast records.

Case studies on synthetic and field datasets, including the Suva Vrela
and Ogun State quarries, show that the proposed method achieves
lower prediction errors than USBM-type empirical formulas
and standard machine learning baselines in the few-shot regime.
Moreover, the digital twin produces smooth, monotonic attenuation curves
and yields equivalent USBM parameters that are compatible with existing standards,
facilitating adoption by practitioners.

The results demonstrate that transferring operator-level physics,
rather than merely fitting data-driven models,
is a promising path to breaking the small-data bottleneck
in blast vibration prediction and to building trustworthy digital twins
for construction and mining projects.

% ========================================
% Acknowledgements
% ========================================
\section*{Acknowledgements}

The authors gratefully acknowledge the support of [funding agencies],
and thank [companies or institutions] for providing data and technical insights.


\bibliographystyle{model1-num-names} % 顺序编码制，适合 AutoCon
\bibliography{refs}


\end{document}
